% =====================================================================
%  Resumo e abstract
% =====================================================================

\begin{resumo}
Quem cultiva plantas em casa decide quando regar ou mudar um vaso de
lugar pela aparência da planta, e só percebe o problema quando ela já
murchou. O FitoSfera propõe trocar essa avaliação a olho por medições do
ambiente, comparadas com o que cada espécie exige. Este relatório
documenta o segundo ciclo do projeto: a ligação do nó sensor a quem
guarda o dado e a formação do conjunto de dados. O nó usa um ESP32 com um
DHT11, para temperatura e umidade do ar, e um sensor de luz BH1750, que
mede até cerca de 120 mil lux. A cada cinco minutos, ele envia um pacote
de 19 bytes por Bluetooth Low Energy, com número de sequência, um
identificador que muda a cada reinício e a luz acumulada desde que ligou,
de modo que uma desconexão não apaga a luz recebida. A equipe escolheu
Bluetooth Low Energy e não WiFi porque o central do produto é o celular
do usuário, que já está na casa. Um gravador escrito pela equipe recebeu
os pacotes durante 24 horas de campanha, gravou cada um antes de
confirmá-lo e produziu um conjunto de dados com 97,6\,\% de
completude. Dezessete requisitos foram verificados por dezoito casos de
teste, e 15 foram atendidos, 1 em parte e 1 ficou sem verificação. O ponto fraco é que o nó não guarda
temperatura e umidade quando o central está longe. No próximo ciclo
entram uma fila local, o pareamento e o aplicativo como central.

\textbf{Palavras-chave}: internet das coisas; ESP32; Bluetooth Low Energy;
monitoramento de plantas; conjunto de dados.
\end{resumo}

\begin{resumo}[Abstract]
\begin{otherlanguage*}{english}
People who grow plants at home decide when to water a plant or move a pot
based on how it looks, and notice problems only after it has wilted.
FitoSfera proposes replacing this judgement by eye with measurements of
the environment, compared with what each species needs. This report
documents the project's second cycle: connecting the sensor node to
whoever stores the data, and building the dataset. The node uses an ESP32
with a DHT11 for air temperature and humidity and a BH1750 light sensor
that measures up to about 120 thousand lux. Every five minutes it sends a
19-byte packet over Bluetooth Low Energy, with a sequence number, an
identifier that changes at every reboot and the light accumulated since
power-up, so that a disconnection does not erase the light received. The
team chose Bluetooth Low Energy over WiFi because the product's central
is the user's phone, which is already in the house. A recorder written by
the team received the packets during 24 hours of data
collection, stored each one before acknowledging it and produced a
dataset with 97.6\,\% completeness. Seventeen requirements were
verified by eighteen test cases, and 15 were met, 1 partially and 1 remained unverified. The weak point
is that the node does not keep temperature and humidity while the central
is away. The next cycle adds a local queue, pairing and the app as the
central.

\textbf{Keywords}: internet of things; ESP32; Bluetooth Low Energy; plant
monitoring; dataset.
\end{otherlanguage*}
\end{resumo}
